\section{复变函数的积分}
\label{复变函数的积分}
\subsection{复变函数积分}
\label{复变函数积分}
\subsubsection{有向曲线}
\label{有向曲线}
设 \(C\) 为平面上给定的一条光滑（或按段光滑）曲线，如果选定 \(C\)
的两个可能方向中的一个作为正方向（或正向），那么我们就把 \(C\)
理解为带有方向的曲线，称为\textbf{有向曲线}。

如果 \(A\) 到 \(B\) 作为曲线 \(C\) 的正向，那么 \(B\) 到 \(A\) 就是曲线
\(C\) 的负向，记为 \(C^{-}\).

\textbf{简单闭曲线的正向}： 简单闭曲线 \(C\) 的正向是指，当曲线上的点 \(P\)
顺此方向前进时，邻近 \(P\) 点的曲线的内部始终位于 \(P\)
点的左方。与之相反的方向就是曲线的负方向。
\subsubsection{积分定义}
\label{积分定义}
设函数 \(w = f(z)\) 定义在区域 \(D\) 内，\(C\) 为区域 \(D\) 内起点为
\(A\) 终点为 \(B\) 的一条光滑的有向曲线。把曲线 \(C\) 任意分成 \(n\)
个弧段，设分点为

\[
A = z_0, z_1, \cdots, z_{k-1}, z_k, \cdots, z_n = B,
\]

在每个弧段 \(\widehat{z_{k-1}z_k}\quad (k = 1, 2, \dots, n )\)
上任意取一点 \(\zeta_k\)。

作和式

\[
S_{n} = \sum_{k=1}^{n} f(\zeta_{k}) \cdot (z_{k} - z_{k-1}) = \sum_{k=1}^{n} f(\zeta_{k}) \cdot \Delta z_{k},
\]

这里 \(\Delta z_{k} = z_{k} - z_{k-1}\)，\(\Delta s_{k}\) 为弧段 \(z_{k-1}z_{k}\) 的长度，记
\(\delta = \max\limits_{1 \leq k \leq n} \{\Delta s_{k}\}\). 当 \(n\) 无限增加且 \(\delta \to 0\) 时，
如果不论对 \(C\) 的分法及 \(\zeta_{k}\) 的取法如何，\(S_{n}\) 有唯一极限，
那么称这极限值为函数 \(f(z)\) 沿曲线 \(C\) 的积分，记为

\[
\int_{C} f(z)dz=\lim_{\delta \to 0}\sum^n_{k=1}f(\zeta_k)\cdot \Delta z_k.
\]
\subsubsection{积分存在的条件和计算法}
\label{积分存在的条件和计算法}
\textbf{存在的条件}： 如果 \(f(z)\) 是连续函数，而 \(C\) 是光滑曲线，则
\(\int_C f(z)dz\)一定存在。

\(\int_C f(z)dz=\int_C udx-vdy+i\int_C vdx+udy\).

\textbf{积分的计算法}： 可以通过两个二元实变函数的线积分来计算。

\begin{align*}
\int_C f(z) \, dz
&= \int_\alpha^\beta \left\{ u[x(t), y(t)] x'(t) - v[x(t), y(t)] y'(t) \right\} \, dt \\
&\quad + i \int_\alpha^\beta \left\{ v[x(t), y(t)] x'(t) + u[x(t), y(t)] y'(t) \right\} \, dt \\
&= \int_\alpha^\beta \left\{ u[x(t), y(t)] + i v[x(t), y(t)] \right\} \left\{ x'(t) + i y'(t) \right\} \, dt \\
&= \boxed{\int_\alpha^\beta f[z(t)] z'(t) \, dt}
\end{align*}

\begin{tcolorbox}[title={Note},colback=red!5!white
,colframe=red!75!black
,colbacktitle=yellow!50!red
,coltitle=red!25!black, fonttitle=\bfseries,subtitle style={boxrule=0.4pt, colback=yellow!50!red!25!white}]
就是把参数方程的复数形式直接带进去。

于是这变成了一个实参数 \(t\) 的定积分，唯一的不同就是被积函数是复变的了。

\end{tcolorbox}

如果 \(C\) 是 \(C_1,C_2,\cdots,C_n\) 等光滑曲线连接成的按段光滑曲线，则

\[
\int_{C} f(z)\mathrm{d} z=\int_{C_{1}} f(z)\mathrm{d} z+\int_{C_{2}} f(z)\mathrm{d} z+\cdots+\int_{C_{n}} f(z)\mathrm{d} z.
\]
\subsubsection{积分的性质}
\label{积分的性质}
复积分与实变函数的定积分有类似的性质。

\begin{enumerate}
\item \textbf{反向路径性质}
\[
   \int_{C} f(z)dz=-\int_{C^{-}}f(z)dz
   \]
\item \textbf{常数倍性质}
\[
   \int_{C} kf(z)dz=k\int_{C} f(z)dz \quad (k 为常数)
   \]
\item \textbf{线性性质}
\[
   \int_{C}[f(z)\pm g(z)]dz=\int_{C}f(z)dz\pm \int_{C}g(z)dz
   \]
\item \textbf{估值定理} 设曲线 \(C\) 的长度为 \(L\), 函数 \(f(z)\) 在 \(C\) 上满足
\(\left |f(z)\right |\leq M\), 那么
\[
   \left|\int_{C}f(z)dz\right|\leq\int_{C}\left |f(z)\right | ds\leq ML
   \]
\end{enumerate}
\subsubsection{柯西-古萨基本定理}
\label{柯西-古萨基本定理}
如果一个复变函数 \(f(z)\) 在一个\textbf{单连通区域} \(D\) 内是\textbf{解析}的，那么 \(f(z)\) 沿 \(D\) 内任意一条\textbf{简单闭合曲线} \(C\) 的积分都等于零。

\begin{tcolorbox}[title={Note},colback=red!5!white
,colframe=red!75!black
,colbacktitle=yellow!50!red
,coltitle=red!25!black, fonttitle=\bfseries,subtitle style={boxrule=0.4pt, colback=yellow!50!red!25!white}]
\begin{enumerate}
\item 这里的 \(C\) 可以不是简单曲线。
\item 这个定理也成为柯西积分定理。
\end{enumerate}

\end{tcolorbox}
\subsubsection{基本定理的推广}
\label{基本定理的推广}
\paragraph{闭路变形原理}
\label{闭路变形原理}
设 \(C_{1}\) 与 \(C_{2}\) 是两条简单闭曲线，\(C_{2}\) 在 \(C_{1}\) 的内部，
\(f(z)\) 在 \(C_{1}\) 与 \(C_{2}\) 所围的多连通域 \(D\) 内解析，
而在 \(D + C_{1} + C_{2}\) 上连续，则

\[
\int_{C_{1}} f(z) \, dz = \int_{C_{2}} f(z) \, dz.
\]

\begin{tcolorbox}[title={注：},colback=red!5!white
,colframe=red!75!black
,colbacktitle=yellow!50!red
,coltitle=red!25!black, fonttitle=\bfseries,subtitle style={boxrule=0.4pt, colback=yellow!50!red!25!white}]
如果我们把如上两条简单闭曲线 \(C_{1}\) 及 \(C_{2}^{-}\) 看成一条复合闭路 \(\Gamma\)，
而且它的正向为：外面的闭曲线 \(C_{1}\) 按逆时针进行，里面的闭曲线 \(C_{2}\) 按
顺时针进行（即沿 \(\Gamma\) 的正向进行时，\(\Gamma\) 的内部总在 \(\Gamma\) 的左手边），那么

\[
\oint_{\Gamma} f(z) \, dz = 0.
\]

该定理说明，在区域内的一个解析函数沿闭曲线的积分，不因闭曲线在区域内作连续变形而
改变它的值，只要在变形过程中曲线\textbf{不经过}函数 \(f(z)\) \textbf{不解析}的点。

\end{tcolorbox}
\paragraph{复合闭路定理}
\label{复合闭路定理}
设 \(C\) 为多连通域 \(D\) 内的一条简单闭曲线，\(C_1, C_2, \ldots, C_n\) 是
在 \(C\) 内部的简单闭曲线，它们互不包含也互不相交，并且以
\(C, C_1, C_2, \ldots, C_n\) 为边界的区域全含于 \(D\)。如果 \(f(z)\) 在 \(D\) 内
解析，那么

\[
\oint_{C} f(z) \, dz = \sum_{k=1}^{n} \oint_{C_k} f(z) \, dz,
\]

其中 \(C\) 及 \(C_k\) (\(k = 1, 2, \ldots, n\)) 均取正方向； 或

\[
\oint_{\Gamma} f(z) \, dz = 0.
\]

其中 \(\Gamma\) 为由 \(C\) 及 \(C_k\) (\(k = 1, 2, \cdots, n\)) 所组成的复合闭路
（其方向是：\(C\) 按逆时针进行，\(C_k\) (\(k = 1, 2, \cdots, n\)) 按顺时针进行）。
\subsubsection{柯西积分公式}
\label{柯西积分公式}
设 \(f(z)\) 在区域 \(D\) 内处处解析，\(C\) 为 \(D\) 内的任何一条正向简单闭曲线，
它的内部完全含于 \(D\)，\(z_0\) 为 \(C\) 内的任一点，那么

\[
f(z_{0}) = \frac{1}{2\pi i} \oint_{C} \frac{f(z)}{z - z_{0}} \, dz.
\]

\begin{tcolorbox}[title={说明},colback=red!5!white
,colframe=red!75!black
,colbacktitle=yellow!50!red
,coltitle=red!25!black, fonttitle=\bfseries,subtitle style={boxrule=0.4pt, colback=yellow!50!red!25!white}]
\begin{enumerate}
\item 把函数在 \(C\) 内部任一点的值用它在边界上的值表示。（这是解析函数的又一特征）
\item 公式不但提供了计算某些复变函数沿闭路积分的一种方法，而且给出了解析函数的一个积分表达式。（这是研究解析函数的有力工具）
\item 一个解析函数在圆心处的值等于它在圆周上的平均值。即，如果 \(C\) 是圆周 \(z = z_{0} + R \cdot e^{i\theta}\)，则有：
$$
    f(z_{0}) = \frac{1}{2\pi} \int_{0}^{2\pi} f(z_{0} + R \cdot e^{i\theta}) \, d\theta
    $$
\end{enumerate}

\end{tcolorbox}
\subsubsection{高阶导数公式}
\label{高阶导数公式}
设函数 \(f(z)\) 在区域 \(D\) 内解析，\(C\) 为 \(D\) 内任意一条包含点 \(z_0\) 的简单正向闭曲线，且 \(C\) 及其内部完全包含于 \(D\)，则 \(f(z)\)
的各阶导数均存在，且 \(f(z)\) 在点 \(z_0\) 处的 \(n\) 阶导数为：

\[
f^{(n)}(z_0) = \frac{n!}{2\pi i} \oint_{C} \frac{f(z)}{(z - z_0)^{n+1}}  dz \quad (n = 1, 2, \cdots)
\]
\subsection{原函数和不定积分}
\label{原函数和不定积分}
\subsubsection{定理}
\label{定理}
\paragraph{定理一}
\label{定理一}
如果 \(f(z)\) 在单连通域内解析，则\textbf{积分与路径无关}。

\begin{tcolorbox}[title={Note},colback=red!5!white
,colframe=red!75!black
,colbacktitle=yellow!50!red
,coltitle=red!25!black, fonttitle=\bfseries,subtitle style={boxrule=0.4pt, colback=yellow!50!red!25!white}]
解析函数在单连通域内的积分只和起点终点有关，那么我们分别让起点终点为
\(z_0\) 和 \(z\)，固定 \(z_0\) 而移动
\(z\)，就可以在这个单连通域内构造一个单值函数
\[F(z) = \int _{z_0}^z f(\zeta)\,\mathrm{d}\zeta\]
你可以把它当作\textbf{势函数}。

\end{tcolorbox}
\paragraph{定理二}
\label{定理二}
如果 \(f(z)\) 在单连通域 \(B\) 内解析，则
\(F(z)=\int_{z_0}^{z}f(\zeta)d\zeta\) 必为 \(B\) 内一个解析函数，使得
\(F'(z)=f(z)\)。

\begin{tcolorbox}[title={Note},colback=red!5!white
,colframe=red!75!black
,colbacktitle=yellow!50!red
,coltitle=red!25!black, fonttitle=\bfseries,subtitle style={boxrule=0.4pt, colback=yellow!50!red!25!white}]
这和微积分中对变上限积分的求导定理很类似，可以类比。

\end{tcolorbox}
\subsubsection{原函数的定义}
\label{原函数的定义}
如果函数 \(\varphi(z)\) 在区域 \(B\) 内的导数为 \(f(z)\)，即
\(\varphi^{\prime}(z)=f(z)\)，那么称 \(\varphi(z)\) 为 \(f(z)\) 在区域
\(B\) 内的原函数。

显然 \(F(z)=\int_{z_0}^{z}f(\zeta)d\zeta\) 是 \(f(z)\) 的一个原函数。

\begin{tcolorbox}[title={原函数之间的关系},colback=red!5!white
,colframe=red!75!black
,colbacktitle=yellow!50!red
,coltitle=red!25!black, fonttitle=\bfseries,subtitle style={boxrule=0.4pt, colback=yellow!50!red!25!white}]
很显然，和实函数差不多的是，\(f(z)\)
的\textbf{任何两个原函数相差一个常数}。

\end{tcolorbox}
\subsubsection{不定积分的定义}
\label{不定积分的定义}
称 \(f(z)\) 的原函数的一般表达式 \(F(z)+c\)（\(c\) 为任意常数）为
\(f(z)\) 的不定积分，记作

\[
\int f(z)dz=F(z)+c.
\]
\paragraph{定理三（类似于牛顿-莱布尼兹公式）}
\label{定理三类似于牛顿-莱布尼兹公式}
如果函数 \(f(z)\) 在单连通域 \(B\) 内处处解析，\(G(z)\) 为 \(f(z)\)
的一个原函数，那么

\[
\int_{z_0}^{z_1} f(z)dz=G(z_1)-G(z_0)
\]

这里 \(z_0, z_1\) 为域 \(B\) 内的两点。

\begin{tcolorbox}[title={Morera 定理},colback=red!5!white
,colframe=red!75!black
,colbacktitle=yellow!50!red
,coltitle=red!25!black, fonttitle=\bfseries,subtitle style={boxrule=0.4pt, colback=yellow!50!red!25!white}]
若函数 \(f(z)\) 在单连通区域 \(D\) 内连续，且对 \(D\)
内任一条逐段光滑的简单闭曲线 \(C\)，都有

$$
\oint_C f(z) \, dz = 0,
$$

则函数 \(f(z)\) 在区域 \(D\) 内解析。

\end{tcolorbox}
\subsection{解析函数和调和函数}
\label{解析函数和调和函数}
\subsubsection{调和函数}
\label{调和函数}
如果二元实变函数 \(\varphi(x,y)\) 在区域 \(D\)
内具有二阶连续偏导数，并且满足拉普拉斯方程

\[
\frac{\partial^{2}\varphi}{\partial x^{2}} + \frac{\partial^{2}\varphi}{\partial y^{2}} = 0,
\]

那么称 \(\varphi(x,y)\) 为区域 \(D\) 内的\textbf{调和函数}。
\subsubsection{解析函数和调和函数的关系}
\label{解析函数和调和函数的关系}
任何在区域 \(D\) 内解析的函数，它的实部和虚部都是 \(D\) 内的调和函数。
\subsubsection{共轭调和函数}
\label{共轭调和函数}
设 \(u(x,y)\) 为区域 \(D\) 内给定的调和函数，我们把使 \(u+iv\) 在 \(D\)
内构成解析函数的调和函数 \(v(x,y)\) 称为 \(u(x，y)\) 的共轭调和函数。

换句话说，在 \(D\) 内满足方程 \[
\frac{\partial u}{\partial x} = \frac{\partial v}{\partial y} ,\quad \frac{\partial u}{\partial y} = -\frac{\partial v}{\partial x}
\] 的两个调和函数中，\(v\) 称为 \(u\) 的共轭调和函数。

\begin{tcolorbox}[title={Note},colback=red!5!white
,colframe=red!75!black
,colbacktitle=yellow!50!red
,coltitle=red!25!black, fonttitle=\bfseries,subtitle style={boxrule=0.4pt, colback=yellow!50!red!25!white}]
区域 \(D\) 内的解析函数的虚部是实部的共轭调和函数。

\end{tcolorbox}
\subsubsection{偏积分法}
\label{偏积分法}
如果已知一个调和函数
\(u\)，那么就可以利用\textbf{柯西黎曼方程}求得它的共轭调和函数，从而构成一个解析函数
\(u+vi\)。这种方法称为\textbf{偏积分法}。
\subsubsection{不定积分法}
\label{不定积分法}
已知调和函数 \(u(x，y)\) 或
\(v(x，y)\)，用不定积分求解析函数的方法称为不定积分法。

\begin{tcolorbox}[title={具体过程},colback=red!5!white
,colframe=red!75!black
,colbacktitle=yellow!50!red
,coltitle=red!25!black, fonttitle=\bfseries,subtitle style={boxrule=0.4pt, colback=yellow!50!red!25!white}]
解析函数 \(f(z)=u+iv\) 的导函数也是解析函数，而且
\[
f'(z) = u_x + iv_x = u_x - iu_y = v_y + iv_x
\]
那我们把 \(u_x - iu_y\) 和 \(v_y + iv_x\) 用 \(z\) 来表示，那么
\[
f'(z) = u_x - iu_y = U(z), \quad f'(z) = v_y + iv_x = V(z).
\]
这两个式子都可以用来积分，分别能得到：
\[
f(z) = \int U(z)\,\mathrm{d}z+c,\quad f(z) = \int V(z)\,\mathrm{d}z+c.
\]
而这两个式子分别可以用于已知实部 \(u\) 和已知虚部 \(v\) 两种情形。

\end{tcolorbox}
